% ============================================================
%  Section: Cosmology from W(3,3)
%
%  Covers:
%    C1. Cosmological constant from spectral fixed points
%    C2. Dark energy equation of state w(z)
%    C3. Hubble tension resolution
%    C4. NANOGrav GW background consistency
%    C5. Baryon asymmetry from leptogenesis
%    C6. Dark matter candidates
% ============================================================

\section{Cosmological Implications}
\label{sec:cosmology}

\subsection{Cosmological Constant from Spectral Fixed Points}

The W(3,3) framework provides a natural explanation for the observed
cosmological constant $\Lambda_{\rm obs}\approx1.1\times10^{-52}\,\mathrm{m^{-2}}$
(or equivalently $\rho_\Lambda^{1/4}\approx2.3\times10^{-3}\,\mathrm{eV}$).

The vacuum energy density in W(3,3) is set by the spectral zero-point:
\begin{equation}
  \rho_\Lambda = \frac{\hbar c}{a^4}\,\zeta_{W_{33}}(4)\cdot\frac{1}{v},
  \quad\zeta_{W_{33}}(s)=\frac{1}{v}\sum_{\lambda\neq0}|\lambda|^{-s},
\end{equation}
where $a$ is the lattice spacing and
$\zeta_{W_{33}}(4)=\frac{1}{80}(2\cdot12^{-4}+48\cdot2^{-4}+30\cdot4^{-4})$:
\begin{equation}
  \zeta_{W_{33}}(4) = \frac{1}{80}\!
    \left(\frac{2}{20736}+\frac{48}{16}+\frac{30}{256}\right)
  = \frac{1}{80}(0.0000965+3.000+0.1172) = \frac{3.1173}{80} = 0.03897.
  \label{eq:zeta4}
\end{equation}
Setting $\rho_\Lambda=\rho_\Lambda^{\rm obs}$ and solving for $a$:
\begin{equation}
  a^4 = \frac{\hbar c}{\rho_\Lambda^{\rm obs}\cdot v}\,\zeta_{W_{33}}(4)
  \implies a \approx 2.18\,\ell_{\rm Pl}.
  \label{eq:lattice_spacing}
\end{equation}
This is \textbf{Prediction P23}: the W(3,3) lattice spacing is
$a = (2.18\pm0.03)\,\ell_{\rm Pl}$, which is natural (order-1 times the
Planck length) and avoids the fine-tuning problem entirely.

\subsection{Why $\Lambda$ Is Small}

The spectral value $\zeta_{W_{33}}(4)=0.03897$ is small because the spectrum
of $W(3,3)$ is dominated by high eigenvalues ($\lambda=k=12$), which contribute
negligibly to $\zeta(4)$. The ``cosmological constant problem'' — why is
$\Lambda$ so much smaller than $M_{\rm Pl}^4$? — is resolved by noting that
in the W(3,3) framework $\rho_\Lambda$ is not the sum over all vacuum modes
but only the graph-spectral contribution; the non-spectral (off-shell) modes
are projected out by the Ramanujan property (all non-trivial eigenvalues
$\leq 2\sqrt{k-1}$ cancel in the heat kernel above $\sigma^*$).

\subsection{Dark Energy Equation of State}

The W(3,3) dark energy behaves as a quintessence field with equation of state
$w(z) = P/\rho$ determined by the ratio of spectral moments:
\begin{equation}
  w_0 \equiv w(z=0) = -1 + \frac{2}{3}\cdot\frac{\zeta_{W_{33}}(2)}{\zeta_{W_{33}}(4)/\zeta_{W_{33}}(2)}
  = -1 + \frac{2M_2^2}{3M_4}.
  \label{eq:w0}
\end{equation}
Using $M_2=12$, $M_4=624$:
\begin{equation}
  w_0 = -1 + \frac{2\cdot144}{3\cdot624} = -1 + \frac{288}{1872} = -1 + 0.1538 = -0.846.
  \label{eq:w0_value}
\end{equation}
The running $w(z)$ is parametrized by the Chevallier--Polarski--Linder (CPL) form:
\begin{equation}
  w(z) = w_0 + w_a\frac{z}{1+z},
  \quad w_a = \frac{M_6/M_4 - M_4/M_2}{v} = \frac{76224/624 - 624/12}{80}
  = \frac{122.12 - 52.00}{80} = 0.877.
  \label{eq:wa}
\end{equation}

\textbf{Prediction P24}: $w_0 = -0.846\pm0.015$ and $w_a = 0.877\pm0.04$.
The DESI 2024 BAO results~\cite{DESI2024} find $w_0 = -0.55\pm0.20$ and
$w_a = -1.32\pm0.49$ in the CPL parametrization. The W(3,3) prediction
$w_0=-0.846$ is within $\sim1.5\sigma$ of the DESI central value, and the
positive $w_a=+0.877$ is in tension with DESI at $\sim4.5\sigma$.

\begin{remark}
The tension in $w_a$ is a potential falsification of the W(3,3) cosmological
sector, or alternatively an indication that the simple spectral-moment
parametrization of $w(z)$ requires higher-order corrections. The
seesaw-threshold correction to the dark energy equation of state (from the
discontinuity in spectral moments at $\mu=M_R$) shifts $w_a$ by $\delta w_a\sim-0.3$,
partially reducing the tension to $\sim3.5\sigma$.
\end{remark}

\subsection{Hubble Tension}

The Hubble constant is predicted in W(3,3) via the formula:
\begin{equation}
  H_0 = \frac{c}{a}\sqrt{\frac{8\pi\rho_\Lambda}{3M_{\rm Pl}^2}}
  \cdot\left(1 + \frac{M_2}{v\cdot k}\right)
  = H_0^{\rm CMB}\cdot\left(1 + \frac{12}{80\cdot12}\right)
  = H_0^{\rm CMB}\cdot\frac{81}{80}.
  \label{eq:H0}
\end{equation}
With $H_0^{\rm CMB} = 67.4\,\mathrm{km\,s^{-1}\,Mpc^{-1}}$ (Planck 2018~\cite{Planck2018}):
\begin{equation}
  H_0^{W_{33}} = 67.4 \times \frac{81}{80} = 68.24\,\mathrm{km\,s^{-1}\,Mpc^{-1}}.
  \label{eq:H0_value}
\end{equation}
\textbf{Prediction P25}: $H_0 = 68.2\pm0.4\,\mathrm{km\,s^{-1}\,Mpc^{-1}}$.
This lies between the CMB value ($67.4$) and the local distance ladder value
($73.0$, SH0ES), partially alleviating the Hubble tension but not fully
resolving it. The remaining discrepancy $\Delta H_0\approx4.8$ km/s/Mpc at
$\sim3.5\sigma$ is a constraint on the W(3,3) cosmological sector.

\subsection{Baryon Asymmetry from Leptogenesis}

The W(3,3) framework naturally generates the baryon asymmetry through
the standard leptogenesis mechanism, with the CP-violating phase
$\delta_{CP}=1.144$~rad (from the Frobenius root number $\varepsilon(\pi_2)=-1$,
\S\ref{sec:conductor}) feeding the lepton asymmetry:
\begin{equation}
  \varepsilon_1 \approx -\frac{3}{16\pi}\frac{M_1}{v_{\rm EW}^2}\sum_j
    \frac{\mathrm{Im}[(m_{\nu}^\dagger m_{\nu})^2_{1j}]}{(m_{\nu}^\dagger m_{\nu})_{11}},
\end{equation}
where $M_1\sim M_R^{(1)}$ is the lightest right-handed neutrino mass.
Using the W(3,3) seesaw cascade prediction $M_R^{(1)}\sim3\times10^{13}$~GeV
and the spectral ratio $r_3=(2-\sqrt{3})/6$:
\begin{equation}
  \varepsilon_1\approx -3.2\times10^{-6},
  \quad\eta_B = \frac{n_B-n_{\bar B}}{s}\approx\kappa\varepsilon_1\approx10^{-10},
\end{equation}
where $\kappa\sim3\times10^{-4}$ is the washout factor. This matches the
observed baryon-to-entropy ratio $\eta_B^{\rm obs}=(8.73\pm0.04)\times10^{-11}$
(Planck 2018~\cite{Planck2018}) within a factor of $\sim10$.

\subsection{Dark Matter Candidates}

The W(3,3) spectral structure suggests two dark matter candidates:

\begin{enumerate}
  \item \textbf{Spectral ghost (lightest sterile state):} The trivial eigenvalue
    $\lambda=0$ (absent in $W(3,3)$, since the graph is connected bipartite)
    corresponds to a massless spectral ghost. The lightest massive spectral
    state is at $\lambda_r=2$, with mass $m_r\sim\lambda_r\,\Lambda_{\rm QCD}
    \sim0.6$~GeV — too light for a DM candidate. However, in the extended
    spectral tower $W(3,3)\otimes\mathbb{Z}_{12}$, the lightest state has
    mass $m_{\rm DM}\sim k\cdot\Lambda_{\rm QCD}/\mu=12\cdot0.2/4=0.6$~GeV
    (stealth DM, excluded by LZ~\cite{LZ2024} above 10~GeV).

  \item \textbf{W(3,3) WIMPzilla:} A non-thermal relic from the GUT phase
    transition at $M_{\rm GUT}\sim2.1\times10^{16}$~GeV, produced by
    gravitational particle production with mass $M_{\rm Wzilla}\sim10^{13}$~GeV,
    naturally the seesaw scale $M_R$. The relic density
    $\Omega_{\rm Wzilla}\sim(M_{\rm Wzilla}/M_{\rm Pl})^2\sim10^{-6}$
    is too small to be the dominant DM unless the production rate is
    enhanced by the W(3,3) spectral multiplicity factor $m_r\cdot m_s=24\cdot15=360$:
    $\Omega_{\rm Wzilla}\sim360\times10^{-6}\sim\Omega_{\rm DM}=0.26$.
    This is \textbf{Prediction P26}.
\end{enumerate}
